\section{Technicians: five methods in one notation}
\label{sec:technicians}

This section restates the technical core of the five readings in a common
notation, sketches the arguments that matter, and ends with the shared
anatomy that makes the works comparable. Throughout, a batch sample is
$x_1,\dots,x_n$, assumed exchangeable unless stated otherwise, with sample
mean $\bar{x}$, sample standard deviation $s$, median $\tilde{x}$ and median
absolute deviation
$\mathrm{MAD} = \operatorname{median}_i \lvert x_i - \tilde{x} \rvert$. A
time series is $x_t$, $t = 1,\dots,N$, observed at a fixed interval $\Delta$.
A detector is a decision rule
$\delta(x) \in \{\text{accept}, \text{flag}\}$, and every detector in this
lineage decomposes into the same three parts: a statistic that measures
surprise, a reference set that defines what unsurprising data look like, and
a threshold that converts surprise into a verdict.

\subsection{Grubbs: outlier detection as hypothesis testing}
\label{sec:tech-grubbs}

\citet{grubbs1969} treats the doubtful observation as a hypothesis test under
the normal model. For a single suspect on either side, the statistic is the
extreme studentised deviate,
\begin{equation}
  G \;=\; \frac{\max_i \lvert x_i - \bar{x} \rvert}{s},
  \label{eq:grubbs}
\end{equation}
and the null hypothesis is that $x_1,\dots,x_n$ are independent draws from
one normal distribution with unknown mean and variance. The elegant part,
due to the line of work Grubbs consolidates, is that the null distribution of
$G$ is available in closed form through a change of viewpoint: conditioning
on which observation is extreme turns the problem into a two-sample $t$
comparison of the suspect against the remaining $n-1$ points, and a
Bonferroni argument over the $n$ possible suspects is exact at the relevant
tail. The two-sided critical value at level $\alpha$ is
\begin{equation}
  G_{\alpha,n} \;=\; \frac{n-1}{\sqrt{n}}
  \sqrt{\frac{t^2_{\alpha/(2n),\,n-2}}{\,n-2+t^2_{\alpha/(2n),\,n-2}\,}},
  \label{eq:grubbscrit}
\end{equation}
where $t_{p,\nu}$ is the upper $p$ quantile of Student's $t$ with $\nu$
degrees of freedom. Section~\ref{sec:experimentalists} verifies this formula
by simulation. Two design decisions in the paper deserve emphasis. First,
Grubbs insists that the analyst declare the alternative in advance, one
outlier or two, upper side or both, because each alternative has its own
statistic; for a pair of suspects on the same side he provides ratio
statistics of the form $S^2_{n-1,n}/S^2$, comparing the residual sum of
squares with and without the suspect pair. Second, he is explicit that the
tests assume a single homogeneous batch. Nothing in the paper licenses
applying Equation~\eqref{eq:grubbs} to a stream with trend or seasonality,
a warning the sensor era would quietly ignore.

The known internal weakness is masking, analysed by \citet{pearson1936}
before the tests were standardised. Two similar outliers inflate $s$
together, so the statistic of each falls below the critical value and the
one-at-a-time test finds nothing. The repair that practice eventually adopted
is Rosner's generalised extreme studentised deviate procedure
\citep{rosner1983}: fix an upper bound $r$ on the number of outliers, remove
the most extreme point and recompute the statistic $r$ times, obtaining
$R_1,\dots,R_r$ with per-step critical values $\lambda_1,\dots,\lambda_r$,
and declare the number of outliers to be the largest $i$ with
$R_i > \lambda_i$. Deciding after all $r$ statistics are computed, rather
than stopping at the first insignificant one, is the whole trick: a masked
pair produces an insignificant $R_1$ but a large $R_2$, which the sequential
stopping rule would never see.

\subsection{Iglewicz and Hoaglin: the robust reformulation}
\label{sec:tech-ih}

\citet{iglewicz1993} replace the fragile ingredients of
Equation~\eqref{eq:grubbs} rather than the logic. The mean has breakdown
point zero, one bad observation moves it arbitrarily, and $s$ is worse,
because contamination inflates it quadratically; the median and the MAD both
withstand corruption of just under half the sample \citep{hampel1974}. Their
modified z-score is
\begin{equation}
  M_i \;=\; \frac{0.6745\,(x_i - \tilde{x})}{\mathrm{MAD}},
  \label{eq:modz}
\end{equation}
with the recommendation to label $x_i$ an outlier when
$\lvert M_i \rvert > 3.5$. The constant is a consistency correction: for
normal data the MAD estimates $\Phi^{-1}(0.75)\,\sigma \approx 0.6745\,
\sigma$, so the factor rescales the MAD into standard deviation units and
makes $M_i$ comparable to a classical z-score. The threshold 3.5 is not a
quantile from a null distribution; the authors calibrated it by simulation on
batches of independent observations as a value that rarely fires on clean
normal samples of moderate size while catching gross errors. The price of
robustness is efficiency, the MAD wastes information relative to $s$ on
genuinely clean data, and the booklet is candid that the rule is a screening
device rather than a test with a controlled error rate. Two properties will
matter later: the scale in the denominator is estimated from whatever set the
analyst supplies, and nothing in the derivation says that set may be a short
sliding window over an autocorrelated stream.

\subsection{Shafer and colleagues: quality assurance as architecture}
\label{sec:tech-shafer}

\citet{shafer2000} contains almost no probability, and that absence is its
content. The Oklahoma Mesonet's quality assurance is an engineered system in
which automated tests are one component among four, alongside laboratory
calibration before deployment, scheduled sensor rotation, and field
intercomparison against travelling reference standards. The automated battery
runs daily over each series and issues, per observation, a graded flag:
0 for good, 1 for suspect, 2 for warning, 3 for failure. In the notation
above, the tests are:
\begin{align}
  \text{range:}\quad & L_p \le x_t \le U_p, \label{eq:range}\\
  \text{step:}\quad & \lvert x_t - x_{t-1} \rvert \le \delta_p, \label{eq:step}\\
  \text{persistence:}\quad & \operatorname{sd}\{x_u : u \in W_t\} \ge
    s^{\min}_p \quad\text{and}\quad
    \max_{u \in W_t} x_u - \min_{u \in W_t} x_u \ge r^{\min}_p,
    \label{eq:persistence}\\
  \text{like-instrument:}\quad & \lvert x^{(1)}_t - x^{(2)}_t \rvert \le d_p,
    \label{eq:like}
\end{align}
where $p$ indexes the parameter, $W_t$ is a trailing window of roughly a
day, and the constants come from instrument specifications and Oklahoma
climatology rather than from any null distribution. The spatial test is the
statistical heart. For station $s_0$ at time $t$, an objective analysis in
the style of \citet{barnes1964} produces a neighbour-based estimate
\begin{equation}
  \hat{x}_t(s_0) \;=\;
  \frac{\sum_{i} w_i\, x_t(s_i)}{\sum_{i} w_i},
  \qquad w_i = \exp\!\big(-r_i^2 / 4\kappa\big),
  \label{eq:barnes}
\end{equation}
with $r_i$ the distance to neighbour $s_i$ and $\kappa$ a smoothing length;
the observation is flagged when $\lvert x_t(s_0) - \hat{x}_t(s_0)\rvert$
exceeds a parameter-specific band. Two architectural principles complete the
design. Flags are advisory: an automated verdict never deletes a value, it
routes the value to a human quality assurance meteorologist whose daily
review, trouble tickets and site visits close the loop. And every verdict is
recorded, so the archive carries both the measurement and the history of
opinions about it. The statistical modesty is deliberate; the system's power
comes from redundancy of evidence, laboratory, field, temporal, spatial,
rather than from the sharpness of any single test.

\subsection{Campbell and colleagues: the streaming framework}
\label{sec:tech-campbell}

\citet{campbell2013} contribute a conceptual framework rather than new
statistics, and the framework is best stated as definitions. Quality
assurance is everything designed before data flow, site selection,
calibration schedules, redundancy, metadata, training; quality control is
everything applied to the data afterwards, automated checks and human
review. The paper's taxonomy separates sensor faults, calibration drift,
fouling, electrical noise, communication dropouts, from anomalies of
scientific interest, real extremes that a naive filter would remove, and
insists the two must be distinguishable in the archive. Its pipeline runs
from sensor through datalogger, telemetry, ingest, automated QC, human
review, to a versioned archive, with three commitments that have become
community norms: never silently delete or overwrite, carry flags as data
alongside the values, and publish data in graded levels, raw, provisional
and quality-controlled, so users choose their own trade-off between latency
and trust. The automated checks themselves are recognisably the Mesonet
battery of Equations~\eqref{eq:range} to~\eqref{eq:barnes}, extended with
ecological examples, and the paper is frank that automation exists to focus
scarce human attention, not to replace it. In the anatomy of this report,
Campbell and colleagues moved the reference set again: from the network's
climatology to the institution's accumulated, versioned record of what its
own sensors do when healthy and when failing.

\subsection{Leigh and colleagues: anomaly detection formalised}
\label{sec:tech-leigh}

\citet{leigh2019} return to explicit statistics, now with the vocabulary of
anomaly detection. Working with turbidity, conductivity and river level
series from in situ sensors in two Queensland catchments, they first define
an anomaly typology grounded in observed failure modes, isolated spikes,
level shifts, persistent drift, periods of anomalously high or low
variability, impossible values, and outages, and have experts label the
historical record accordingly. Three families of detectors are then
compared on the labelled data.

Rule-based methods are the inherited battery: bounds and step limits as in
Equations~\eqref{eq:range} and~\eqref{eq:step}, which catch impossible
values and outages cheaply. Regression-based methods fit a time series model
to a training window and flag observations that fall outside one-step-ahead
prediction intervals: with an ARIMA or regression-with-ARIMA-errors model
producing forecast $\hat{x}_t$ and forecast standard error
$\hat{\sigma}_t$, the rule is
\begin{equation}
  \text{flag } x_t \iff
  \lvert x_t - \hat{x}_t \rvert > k\,\hat{\sigma}_t,
  \label{eq:regression}
\end{equation}
optionally run both forwards and backwards through the series. Feature-based
methods, developed in the companion paper \citep{talagala2019}, embed
windows of the series into a feature space, spectral entropy, autocorrelation
structure, variance ratios and similar summaries, and declare anomalies where
the feature vector is far from the mass of the data; the distance threshold
is not hand-set but derived from extreme value theory, fitting a generalised
Pareto tail to nearest-neighbour distances so that the flagging rate has an
interpretable exceedance probability. This is the lineage's first
principled return, since Grubbs, to thresholds with a probability attached.

The evaluation is the paper's most consequential contribution. Performance
is reported per anomaly type, and no family dominates: regression-based
detectors are strong on the high-priority sudden faults, spikes and shifts,
but degrade when the water itself behaves unusually; feature-based detectors
are conservative, high specificity at the cost of sensitivity; rules are
unbeatable on impossible values and outages; and slow drift defeats
everything that looks only at one-step surprise. The recommendation is an
ensemble with union and intersection rules tuned to the cost structure, and
the framework closes with the same institutional move as Shafer and
Campbell: detectors propose, people dispose.

\subsection{The shared anatomy}
\label{sec:tech-anatomy}

Table~\ref{tab:anatomy} aligns the five works on the three-part anatomy.
Read down the reference set column: theory (a parametric null), then robust
theory (the sample's own resistant core), then engineering and climatology
(what instruments and Oklahoma allow), then institutional memory (versioned
archives of a network's own behaviour), then learned structure (models and
feature distributions fitted to labelled history). That migration is the
technical content of 170 years of this literature compressed into one
sentence: the mathematics of the surprise statistic changed remarkably
little, while the definition of unsurprising data absorbed ever more
context. The thresholds, meanwhile, changed provenance five times without
ever becoming self-justifying, a theme the Experimentalists' and Critics'
sections take up quantitatively.

\begin{table}[t]
  \centering
  \small
  \caption{The three-part anatomy across the reading list.}
  \label{tab:anatomy}
  \begin{tabular}{@{}p{2.9cm}p{3.4cm}p{3.6cm}p{4.4cm}@{}}
    \toprule
    Work & Surprise statistic & Reference set & Threshold provenance \\
    \midrule
    \citet{grubbs1969} & extreme studentised deviate
      & i.i.d.\ normal batch & exact null distribution, chosen $\alpha$ \\
    \citet{iglewicz1993} & modified z-score
      & robust core of the batch & simulation on clean batches; 3.5 \\
    \citet{shafer2000} & rule residuals; deviation from Barnes estimate
      & instrument specs, climatology, neighbours & engineering constants,
        tuned in operations \\
    \citet{campbell2013} & inherited battery
      & institutional archive of healthy behaviour & site policy, versioned
        and reviewed \\
    \citet{leigh2019} & forecast innovation; feature-space distance
      & fitted model of labelled history & prediction intervals; extreme
        value tail on distances \\
    \bottomrule
  \end{tabular}
\end{table}
